problem:  
Let \( (M^n,g) \) be a compact, simply connected Riemannian manifold with positive sectional curvature.  
At each \(p\in M\), define the pointwise sectional curvature pinching ratio  
\[
\delta(p)=\frac{\min_{\sigma\subset T_pM}K_p(\sigma)}{\max_{\sigma\subset T_pM}K_p(\sigma)}.
\]  
Assume there is a constant \(0<\varepsilon\ll1\) such that the *almost‑quarter‑pinched region*
\[
E_\varepsilon=\{p\in M:\tfrac{1}{4}-\varepsilon\le\delta(p)\le\tfrac{1}{4}+\varepsilon\}
\]
satisfies \(\operatorname{Vol}(E_\varepsilon)\ge c\,\operatorname{Vol}(M)\) for some fixed \(c>0\) independent of \(\varepsilon\).  

**Problem:**  
Does the existence of such a region \(E_\varepsilon\) of uniform positive volume fraction for sufficiently small \(\varepsilon\) force \(M\) to be diffeomorphic to a spherical space form?  
Equivalently, can there exist a sequence of positively curved, simply connected manifolds \((M_i^n,g_i)\) with  
\[
\operatorname{Vol}(E_{\varepsilon_i}(M_i))/\operatorname{Vol}(M_i)\ge c>0,\qquad \varepsilon_i\to0,
\]  
but with \(M_i\) not diffeomorphic to \(S^n\)?  
If so, what is the limit geometry (in the Cheeger–Gromov sense) of such nearly quarter‑pinched metrics: does curvature concentration or collapse inevitably occur along the almost‑critical set \(E_{\varepsilon_i}\)?  


Why is it a "good" problem:  
This problem cleanly isolates the **stability frontier of spherical rigidity** under the \(1/4\)-pinching condition, in a quantitative, measurable form. Classical sphere theorems require strict pinching everywhere; this setting relaxes that to a *substantial region of near-critical curvature*, asking whether topology already rigidifies before global uniformity is reached. It inquires into the geometric threshold at which *macroscopic quarter‑pinching* enforces global spherical topology.

A resolution would illuminate how **local curvature uniformity influences global topology**, deepening the interplay between comparison geometry, Ricci flow dynamics, and the topology of positively curved manifolds. The volume‑fraction condition sharpens the measure‑theoretic aspect, introducing a natural quantitative setting akin to \(\varepsilon\)-regularity theory in geometric analysis.  

Its formulation is simple—requiring only curvature ratios and volume proportions—but its resolution would require a synthesis of analytic (Ricci flow stability), geometric (Cheeger–Gromov convergence), and topological (rigidity classification) insights. Hence, it stands as a crisp, open, and deeply integrative question at the front line of Riemannian geometry.